%% Week 43 -- Tuesday session (flipped classroom): first the worked screening
%% example of Monday, measure by measure (Part A, ~25 min in plenum); then
%% ONE case (Part B): the confusion matrix, the ROC curve and the AUC written
%% from the definitions and checked against scikit-learn, on "is it an
%% eight?" with the scikit-learn digits and logistic regression against a
%% one-hidden-layer network -- the structure of exercises 1 and 2 of
%% exercisesweek43.ipynb (deadline Friday October 23).  No validation set and
%% no threshold selection: cross-validation for the Project 2 network is the
%% Tuesday session and the exercises of week 44.  The case is followed by a
%% "do you have the insight?" frame (answers on \pause).
%% Companion to week43slides.tex (Monday lecture) and to week43tuesday.ipynb.
%% Sources: Section 5.12 of the book, Section 2.4 for the correlation
%% coefficient, the week-41 network code; all numbers from
%% week43tuesday.ipynb (seed 2026).
%% Build:  pdflatex week43tuesday.tex  (twice)
\documentclass[10pt,aspectratio=169]{beamer}
\usetheme{red_plain}
\usepackage[T1]{fontenc}
\usepackage[utf8]{inputenc}
\usepackage{amsmath,amssymb,bm}
\usepackage{graphicx}
\usepackage{listings}
\usepackage{xcolor}
\usepackage{booktabs}
\graphicspath{{../BookFigures/chapter05_logistic_regression/}{./}}

%% ---- code listings
\definecolor{codebg}{rgb}{0.97,0.97,0.97}
\lstset{language=Python,basicstyle=\ttfamily\scriptsize,backgroundcolor=\color{codebg},
  keywordstyle=\color{blue!70!black}\bfseries,commentstyle=\color{green!45!black},
  stringstyle=\color{red!60!black},showstringspaces=false,frame=single,
  framerule=0pt,xleftmargin=4pt,xrightmargin=4pt,breaklines=true,columns=fullflexible}

%% ---- discussion frames (kept for future weeks, unused in this deck)
\newcounter{insight}
\newenvironment{insight}[1][]{%
  \stepcounter{insight}%
  \begin{frame}{Warm-up \arabic{insight}: #1}%
  \setbeamercolor{block title}{bg=structure.fg,fg=white}%
  \setbeamercolor{block body}{bg=structure.fg!8}%
}{\end{frame}}
\newcommand{\question}[1]{\begin{block}{Question}#1\end{block}}
\newcommand{\answer}[1]{\pause\begin{block}{Answer}#1\end{block}}
\newcommand{\hint}[1]{\pause{\small\emph{Hint: }#1}}

%% ---- case frames
\newcounter{case}
\newenvironment{case}[1][]{%
  \stepcounter{case}%
  \begin{frame}[fragile,environment=case]{Case \arabic{case}: #1}%
}{\end{frame}}

\setbeamertemplate{footline}[frame number]
\setbeamertemplate{blocks}[rounded][shadow=false]
\setbeamerfont{block title}{size=\small}
\setbeamercovered{invisible}

\title{Week 43, Tuesday Session}
\subtitle{Flipped classroom: the measures of a classifier, measure by
  measure -- then one case: your own confusion matrix, ROC curve and AUC\\[3pt]
  \small Bring your laptop; open \texttt{week43tuesday.ipynb} and
  \texttt{exercisesweek43.ipynb}}
\author{Morten Hjorth-Jensen}
\institute{Department of Physics and Center for Computing in Science Education,\\
  University of Oslo, Norway}
\date{}

\begin{document}

\begin{frame}
\titlepage
\end{frame}

\begin{frame}{How today works}
\small
\begin{block}{The format (as every Tuesday)}
Mondays are lectures; Tuesdays belong to your questions and to concrete
cases.  Today has two parts, and the second is the skeleton of this week's
exercises and of part e) of Project 2.
\end{block}
\begin{enumerate}
  \item \textbf{Part A, first $\sim$25 min, in plenum:} Monday's worked
        screening example measure by measure -- the four counts, accuracy
        and why it prefers the useless test, precision and recall and their
        two denominators, $F_1$, balanced accuracy, MCC, the AUC as a
        probability about pairs, and why we train on the cross entropy.
        Bring your questions on Monday's lecture here.
  \item \textbf{Part B, the rest of the double session, in groups:} Case 1
        with \texttt{week43tuesday.ipynb}: your own confusion matrix, ROC
        curve and AUC, written from the definitions and checked against
        \texttt{scikit-learn}, for logistic regression and a one-hidden-layer
        network from your own code on ``is it an eight?'' -- then the
        \texttt{scikit-learn} calls you will use in Project 2, and the
        insight questions in plenum.
  \item \textbf{Last 5 min:} wrap-up, the exercises, and next week.
\end{enumerate}
\begin{block}{Exercises week 43}
Exercise 1: the functions of Part B, with a small test.  Exercise 2: MNIST,
one digit against the rest, with your Project 2 classifiers.  Exercise 3
(optional): the $10\times10$ matrix.  Deadline \textbf{Friday October 23 at
midnight}; declare your use of LLMs.  No validation sets yet -- that is
week 44.
\end{block}
\end{frame}

%% ====================================================================
%% Part A: the worked example
%% ====================================================================

\begin{frame}{Part A: the worked example -- four counts and seven numbers}
\small
\begin{columns}[T]
\begin{column}{0.46\textwidth}
\begin{block}{A screening test, prevalence $\pi=2\%$, $1000$ people}
\scriptsize
\[
  \begin{array}{c|cc}
     & \text{predicted } 1 & \text{predicted } 0\\
    \hline
    \text{true } 1 & \mathrm{TP}=15 & \mathrm{FN}=5\\
    \text{true } 0 & \mathrm{FP}=40 & \mathrm{TN}=940
  \end{array}
\]
Twenty people have the condition; the test finds fifteen and raises forty
false alarms.  Rows are the truth, columns the decision; FN are the
\emph{missed} cases (the dangerous error in screening), FP the \emph{false
alarms} (a follow-up, not a life).  Next to it: the ``test'' that declares
everybody healthy, $0,20,0,980$.
\end{block}
\begin{block}{Ask the room first}
\scriptsize
Before showing the table: which of the two tests has the higher accuracy?
Which would you rather be screened with?  Why do the two answers differ?
\end{block}
\end{column}
\begin{column}{0.52\textwidth}
\begin{center}
\scriptsize
\setlength{\tabcolsep}{3pt}
\begin{tabular}{@{}llcc@{}}
\toprule
measure & definition & test & empty\\
\midrule
accuracy & $(\mathrm{TP}+\mathrm{TN})/n$ & $0.955$ & $\mathbf{0.980}$\\
precision & $\mathrm{TP}/(\mathrm{TP}+\mathrm{FP})$ & $0.273$ & $0$ ($0/0$)\\
recall, TPR & $\mathrm{TP}/(\mathrm{TP}+\mathrm{FN})$ & $0.750$ & $0$\\
specificity, TNR & $\mathrm{TN}/(\mathrm{TN}+\mathrm{FP})$ & $0.959$ & $1.000$\\
$F_1$ & $2\,\mathrm{prec}\cdot\mathrm{rec}/(\mathrm{prec}+\mathrm{rec})$ & $0.400$ & $0$\\
balanced acc. & $\tfrac12(\mathrm{TPR}+\mathrm{TNR})$ & $0.855$ & $0.500$\\
MCC & $(\mathrm{TP}\,\mathrm{TN}-\mathrm{FP}\,\mathrm{FN})/\sqrt{D}$ & $0.436$ & $0$\\
\bottomrule
\end{tabular}
\end{center}
\scriptsize
$D=(\mathrm{TP}+\mathrm{FP})(\mathrm{TP}+\mathrm{FN})(\mathrm{TN}+\mathrm{FP})(\mathrm{TN}+\mathrm{FN})$.
What each row asks: decisions right, all alike; of those flagged, how many
were?; of the sick, how many found?; of the healthy, how many left alone?;
both of the previous two; accuracy at equal class weights; correlation of
$y$ and $\tilde y$.  Accuracy ranks the two tests the wrong way round;
every other row ranks them the right way, and each says something
different.
\end{column}
\end{columns}
\end{frame}

\begin{frame}{Part A: what the numbers mean, one at a time}
\scriptsize
\begin{columns}[T]
\begin{column}{0.5\textwidth}
\begin{block}{Accuracy is a prevalence-weighted average}
With $n_+=\mathrm{TP}+\mathrm{FN}$, $n_-=\mathrm{FP}+\mathrm{TN}$,
$\pi=n_+/n$:
\[
  \mathrm{accuracy}=\frac{n_+}{n}\frac{\mathrm{TP}}{n_+}+\frac{n_-}{n}\frac{\mathrm{TN}}{n_-}
  =\pi\,\mathrm{recall}+(1-\pi)\,\mathrm{specificity}
\]
$=0.02\times0.75+0.98\times0.959=0.955$: the sick enter with weight
$0.02$.  Compare every accuracy with the majority rule $\max(\pi,1-\pi)$
and quote the fraction of the \emph{remaining} error removed,
$1-(1-\mathrm{acc})/(1-p_0)$ -- here negative.  Balanced accuracy is the
same average at equal weights: $0.855$ against $0.5$ for a coin.
\end{block}
\begin{block}{Recall is about the misses, precision about the false alarms}
Different denominators: the truth ($n_+=20$) for recall, our decision
($\mathrm{TP}+\mathrm{FP}=55$) for precision.  Three quarters of the sick
are found; three of four people called back are healthy, because the
condition is rare ($40$ false alarms among $980$ healthy is a $4\%$ false
positive rate).  Lower the threshold: recall up, precision down.  Screening
wants recall, a spam filter precision.
\end{block}
\end{column}
\begin{column}{0.48\textwidth}
\begin{block}{$F_1$ is small whenever either is small}
With $m=\min(\mathrm{prec},\mathrm{rec})$, $M=\max$:
$F_1=2mM/(m+M)=2m/(1+m/M)$, so $m\le F_1\le2m$.  That is why the harmonic
mean: precision $1$ and recall $0.1$ give an arithmetic mean of $0.55$ and
a harmonic mean of $0.18$ -- perfect on one axis and useless on the other
is useless.  Worked example: $2\times0.273\times0.75/1.023=0.40$.  $F_1$
never looks at TN.
\end{block}
\begin{block}{MCC is a correlation}
The Pearson correlation (Section 2.4) between the $0/1$ labels and the
$0/1$ decisions: all four counts, symmetric in the classes, $0$ for the
empty rule, $\pm1$ perfect/perfectly wrong -- hard to fool by imbalance:
$13\,900/31\,917=0.436$.
\end{block}
\begin{block}{Why we train on the cross entropy}
All of the above are functions of the decisions: staircases in the
parameters, zero gradient.  The cross entropy is the smooth surrogate,
minimised \emph{roughly} where the decisions are right: one positive, one
negative, model A ($0.51$, $0.49$) has accuracy $1$ and cross entropy
$0.673$, model B ($0.99$, $0.60$) accuracy $0.5$ and $0.463$.
\end{block}
\end{column}
\end{columns}
\end{frame}

\begin{frame}{Part A: the AUC is a probability about pairs}
\scriptsize
\begin{columns}[T]
\begin{column}{0.5\textwidth}
\begin{block}{From one threshold to all of them}
A confusion matrix is one threshold: one point
$(\mathrm{FPR},\mathrm{TPR})=(0.041,0.75)$ for our test.  Scores give a
whole curve, from $(0,0)$ (flag nobody) to $(1,1)$ (flag everybody); a
perfect ranking passes through $(0,1)$, a random one is the diagonal.  The
curve is every trade-off the model offers; a threshold picks one point.
\end{block}
\begin{block}{The rank formula}
Sort all $n=n_++n_-$ scores increasingly; let $R_i$ be the rank of the
$k_i$th positive.  Below it sit $R_i-1$ samples, $k_i-1$ of them positives,
so it outranks $R_i-k_i$ negatives.  Summed over the positives:
$\sum_iR_i-\tfrac12n_+(n_++1)$ correctly ordered pairs of $n_+n_-$,
\[
  \mathrm{AUC}=\frac{\sum_{i\in\mathrm{pos}}R_i-\tfrac12n_+(n_++1)}{n_+n_-}
  =\mathrm{Prob}(\text{pos}>\text{neg}),
\]
the Mann--Whitney statistic (ties: average rank, half a pair), equal to
the trapezoidal area under the curve.
\end{block}
\end{column}
\begin{column}{0.48\textwidth}
\begin{block}{Three checks of the formula}
A random ranking: every pair is ordered correctly with probability
$\tfrac12$, so $0.5$.  All scores equal: every $R_i=(n+1)/2$ with average
ranks, $\sum_iR_i=n_+(n+1)/2$, numerator $n_+n_-/2$: again $0.5$.  Perfect
separation: the positives hold the top $n_+$ ranks,
$\sum_iR_i=n_+n_-+\tfrac12n_+(n_++1)$: $1$.
\end{block}
\begin{block}{What the AUC ignores}
The threshold -- there is none in it.  The prevalence -- duplicate every
positive: TP and FN double together, every TPR is unchanged, FPR never saw
the positives.  Precision is different: FP sits in its denominator, so the
precision--recall curve moves up when positives are duplicated, and its
floor is $\pi$ -- the curve to look at when positives are rare and the ROC
curve flatters ($1\%$ FPR on $9108$ negatives is $91$ false alarms against
$892$ fives).
\end{block}
\begin{block}{In one sentence each}
Matrix: what a user of the classifier experiences at a stated threshold.
AUC: what the model is capable of.  A report gives both.
\end{block}
\end{column}
\end{columns}
\end{frame}

%% ====================================================================
%% Part B: the case
%% ====================================================================

\begin{case}[your own confusion matrix, ROC curve and AUC, in five steps]
\footnotesize
\textbf{Goal:} leave with your own \texttt{confusion\_matrix},
\texttt{classification\_report}, \texttt{roc\_curve} and \texttt{roc\_auc},
agreeing with \texttt{scikit-learn} to every digit, and with the
\texttt{scikit-learn} calls for Project 2 at your fingertips (notebook,
Case 1; seed 2026).  Exercises 1 and 2 are the same steps.
\begin{enumerate}
  \item \textbf{The data and the number to beat:} the $8\times8$ digits,
        $y=1$ for the eights ($9.7\%$); stratified $80/20$ split, $1437$
        training and $360$ test images.  ``Never an eight'' scores
        $0.9028$ -- write it down before training anything.
  \item \textbf{Two classifiers with your own code:} logistic regression
        (the week-41 network with no hidden layer, $64\to1$, $65$
        parameters) and one hidden layer of $32$ ReLU units ($2113$
        parameters), SGD with $M=32$, $\gamma=0.1$, $150$ epochs.
  \item \textbf{The confusion matrix by hand:} four counts for $p>0.5$,
        then accuracy, precision, recall, specificity, $F_1$, balanced
        accuracy and MCC from the definitions of Part A -- checked against
        \texttt{sklearn.metrics} (mind its layout).
  \item \textbf{The ROC curve and the AUC from one sorted pass:} the curve
        from cumulative TP and FP down the sorted scores, the area by the
        trapezoidal rule and by the rank formula; the histogram of the
        output probabilities by class; all against \texttt{roc\_auc\_score}
        and \texttt{roc\_curve}.
  \item \textbf{The same with \texttt{scikit-learn}:}
        \texttt{confusion\_matrix}, \texttt{classification\_report},
        \texttt{roc\_curve}, \texttt{roc\_auc\_score},
        \texttt{ConfusionMatrixDisplay}, \texttt{RocCurveDisplay} -- the
        template for Project 2.
\end{enumerate}
\end{case}

\begin{frame}[fragile]{Case 1: the code you should end up with}
\begin{columns}[T]
\begin{column}{0.5\textwidth}
\begin{lstlisting}
def confusion_matrix(y, yhat):
    tp = np.sum((y == 1) & (yhat == 1))
    fn = np.sum((y == 1) & (yhat == 0))
    fp = np.sum((y == 0) & (yhat == 1))
    tn = np.sum((y == 0) & (yhat == 0))
    return tp, fn, fp, tn

def roc_curve(y, s):
    order = np.argsort(-s, kind="mergesort")
    y, s = y[order], s[order]
    tp = np.cumsum(y == 1)
    fp = np.cumsum(y == 0)
    last = np.r_[np.diff(s) != 0, True]
    tpr = np.r_[0.0, tp[last] / tp[-1]]
    fpr = np.r_[0.0, fp[last] / fp[-1]]
    return fpr, tpr, np.r_[np.inf, s[last]]

def roc_auc(y, s):            # P(pos > neg)
    _, inv, cnt = np.unique(
        s, return_inverse=True, return_counts=True)
    first = np.cumsum(cnt) - cnt + 1
    ranks = (first + (cnt - 1) / 2)[inv]
    n_pos = np.sum(y == 1)
    n_neg = len(y) - n_pos
    S = np.sum(ranks[y == 1])
    S -= n_pos * (n_pos + 1) / 2
    return S / (n_pos * n_neg)
\end{lstlisting}
\end{column}
\begin{column}{0.48\textwidth}
\begin{lstlisting}
# step 5: the scikit-learn template for Project 2
from sklearn.metrics import (
    confusion_matrix, classification_report,
    precision_score, recall_score, f1_score,
    roc_curve, roc_auc_score,
    ConfusionMatrixDisplay, RocCurveDisplay)

y_pred = p_test > 0.5
cm = confusion_matrix(y_test, y_pred)
print(cm)                 # [[TN, FP], [FN, TP]]
print(classification_report(y_test, y_pred))
fpr, tpr, thr = roc_curve(
    y_test, p_test, drop_intermediate=False)
print(roc_auc_score(y_test, p_test))

fig, ax = plt.subplots(1, 2, figsize=(9, 3.6))
ConfusionMatrixDisplay(cm).plot(ax=ax[0])
RocCurveDisplay.from_predictions(
    y_test, p_test, ax=ax[1])
\end{lstlisting}
\scriptsize
Checkpoints: your ROC curve has $361$ points and equals
\texttt{roc\_curve(..., drop\_intermediate=False)}; the three AUCs agree
to five decimals.
\end{column}
\end{columns}
\end{frame}

\begin{frame}{Case 1: what you should find}
\scriptsize
\begin{columns}[T]
\begin{column}{0.5\textwidth}
\begin{block}{Steps 2--4: the test set at $t=0.5$}
\begin{center}
\setlength{\tabcolsep}{3pt}
\begin{tabular}{@{}lrrrrrrr@{}}
\toprule
 & TP & FN & FP & acc & prec & rec & AUC\\
\midrule
never an eight & $0$ & $35$ & $0$ & $0.903$ & -- & $0$ & $0.5$\\
logistic regr. & $25$ & $10$ & $3$ & $0.964$ & $0.893$ & $0.714$ & $0.9613$\\
one hidden layer & $29$ & $6$ & $2$ & $0.978$ & $0.935$ & $0.829$ & $0.9896$\\
\bottomrule
\end{tabular}
\end{center}
TN $=322$ and $323$ of $325$.  The accuracies are $0.96$ and $0.98$ against
the free $0.90$: the two classifiers removed $63\%$ and $77\%$ of the
remaining error.  Logistic regression misses $10$ of $35$ eights; $F_1$
$0.79$ against $0.88$, balanced accuracy $0.85$ against $0.91$, MCC $0.78$
against $0.87$.  The hand-made numbers agree with \texttt{scikit-learn} to
every digit printed; the hand-made ROC curve has $361$ points and is
identical to \texttt{roc\_curve(..., drop\_intermediate=False)}; the three
AUCs agree to five decimals.
\end{block}
\end{column}
\begin{column}{0.48\textwidth}
\begin{block}{Reading it with Part A}
Accuracy $=0.097\times\mathrm{recall}+0.903\times\mathrm{specificity}$:
for logistic regression $0.097\times0.714+0.903\times0.991=0.964$ -- the
missed eights cost it $0.03$, the three false alarms $0.008$.  The two
classifiers differ almost only in FN ($10$ against $6$): recall is the
measure that moves ($0.71\to0.83$), precision barely
($0.89\to0.94$).
\end{block}
\begin{block}{What the AUC adds}
Logistic regression at $0.5$ misses $10$ eights, yet its AUC is $0.96$:
their $p$ runs from $0.006$ to $0.45$, and each of them is still ranked
above $157$ to $321$ of the $325$ non-eights.  A lower threshold would
find some of them at the price of false alarms -- the operating point can
be slid along the curve; the histogram of step 4 shows what it would cost.
Which point to choose, and on which data, is next week.
\end{block}
\begin{block}{Step 5}
\texttt{classification\_report} prints precision, recall and $F_1$ for
\emph{both} classes plus their macro and weighted averages: $0.94/0.83/0.88$
for the eights, $0.98/0.99/0.99$ for the rest -- the same numbers, with the
negatives' side made explicit.
\end{block}
\end{column}
\end{columns}
\end{frame}

\begin{frame}{Case 1: do you have the insight?}
\scriptsize
\vspace{-3pt}
\begin{enumerate}
  \item \textbf{Accuracy.}  Both classifiers scored above $0.95$ while
        ``never an eight'' scores $0.90$.  What fraction of the
        \emph{remaining} error did each remove?  Write each accuracy as
        $\pi\,\mathrm{recall}+(1-\pi)\,\mathrm{specificity}$ and say which
        term carries it.  Which single numbers survive the imbalance?
  \item \textbf{The two errors.}  The two classifiers differ mostly in one
        of the four counts.  Which one, what does it mean for an eight, and
        which measure is it that moves most between the two?
  \item \textbf{The AUC.}  State in one sentence, about pairs of images,
        what an AUC of $0.99$ means.  Why does the ROC curve ignore the
        threshold and the class balance, and why does the precision--recall
        curve not?
  \item \textbf{What a confusion matrix cannot say.}  The AUC of logistic
        regression is $0.96$, not $0.5$, although it misses ten eights at
        $t=0.5$.  What does that tell you about the choice of threshold, and
        why is the matrix at $0.5$ nevertheless what Project 2 asks for?
\end{enumerate}
\pause
\begin{block}{What we hope you say}
\scriptsize
(1) With $p_0$ the majority accuracy, the fraction removed is
$1-(1-\mathrm{acc})/(1-p_0)$: $63\%$ and $77\%$.  The accuracy is
$0.1\times\mathrm{recall}+0.9\times\mathrm{specificity}$, carried almost
entirely by the specificity, so missing a quarter of the eights still
scores above $0.95$; balanced accuracy, $F_1$ and the MCC give the majority
rule $0.5$ or $0$ and do not reward it.
(2) FN, the missed eights; recall is the measure that moves,
$0.71\to0.83$.
(3) Pick a random eight and a random non-eight: with probability $0.99$
the eight gets the higher score.  The ROC curve is built from the ranking
alone, and copying every eight leaves every rate unchanged; precision has
FP in its denominator, so the fraction of positives enters it.
(4) $0.5$ is a convention; a lower threshold would find more of the eights
at the price of false alarms -- the operating point can be moved along the
curve, and choosing it, on validation data and from the costs of the two
errors, is next week.  The matrix at a stated threshold is what a user of
the classifier experiences; the AUC is what the model is capable of.  A
report gives both.
\end{block}
\end{frame}

\begin{frame}{Wrap-up, the exercises, and next week}
\small
\begin{block}{What you should own now}
\footnotesize
A confusion matrix you can compute and read without a library; accuracy,
precision, recall and $F_1$ from their definitions and from the table of
Part A, with the two identities ($\mathrm{acc}=\pi\,\mathrm{rec}+(1-\pi)\,
\mathrm{spec}$, $m\le F_1\le2m$); the ROC curve and the AUC from a sorted
pass and from the rank formula, checked against \texttt{scikit-learn}; and
the \texttt{scikit-learn} calls for Project 2.
\end{block}
\begin{block}{Exercises week 43 -- deadline Friday October 23 at midnight}
\footnotesize
\texttt{exercisesweek43.ipynb}: exercise 1 is steps 3--5 as functions
with a small test (tied scores on purpose); exercise 2 is steps 1--4 on
MNIST, one digit against the rest, with your Project 2 classifiers;
exercise 3 (optional) is the $10\times10$ matrix.  Hand in one notebook per
group in Canvas; list your collaborators; finish the self-check table;
declare your use of LLMs.
\end{block}
\begin{block}{Project 2 (deadline November 6) and next week}
\footnotesize
Part e) of Project 2 asks for logistic regression as the baseline and for
the confusion matrix on the full MNIST test set with its interpretation:
today's steps 3--5 on the full data, read with Part A.  \textbf{Week 44,
Tuesday session and exercises:} validation sets and cross-validation for
your Project 2 network -- choosing the hyperparameters, and the threshold,
without touching the test set.  \textbf{Week 44, Monday:} Chapter 10
continued -- equivariance and pooling.
\end{block}
\end{frame}

\end{document}
